\documentclass[12pt]{article}
\usepackage[utf8]{inputenc}
\usepackage[T1]{fontenc}
\usepackage{hyperref}
\title{Poqomam, Southern Language Resources for K-12 Curriculum Development}
\author{Compiled by Claude Sonnet 5.5 (claude-sonnet-5-5) for Andr\'as Kornai}
\date{}
\begin{document}
\maketitle

\section{Introduction}
The code \texttt{pou} (Poqomam, Southern) dates from Ethnologue's 2009 edition, which split Poqomam into three languages; the 2013 edition reverted to a single code \texttt{poc}. The Southern varieties (Palín, Escuintla; Santa Cruz Chinautla, Guatemala; San Luis Jilotepeque, Jalapa) are therefore usually documented under the general name Poqomam (Mayan, Guatemala). The ``Bookkeeping'' family label in the request is a Glottolog artefact, not a genuine family. This report is based on one web search and is not exhaustive.

\section{Primary Educational Resources}
No school textbooks or primers were found in this search.

\section{Linguistic and Reference Materials}
\begin{itemize}
\item Academia de Lenguas Mayas de Guatemala, \emph{Diccionario Bilingüe Poqomam-Español} (\emph{Rukorikb'aal tz'ihb'anik / sirinik reh poqom q'orb'al}), covering San Luis Jilotepeque, Chinautla and Palín; reported by AGN, December 2024. \url{https://agn.gt/?p=402251}
\item Guatemalan Mayan Dialect Survey, Poqomam of San Luis Jilotepeque: 1971 wordlist collected by Terrence Kaufman, AILLA. \url{https://islandora-ailla.lib.utexas.edu/islandora/object/ailla%3A137807}
\item Endangered Languages Project entry for Poqomam. \url{https://endangeredlanguages.com/lang/4986}
\end{itemize}

\section{Religious and Community Resources}
\begin{itemize}
\item Catholic Church of Guatemala, Bible resources in Mayan languages (Poqom language document). \url{https://www.iglesiacatolica.org.gt/bibliakiche/BKiche_Lenguas_LenguasPoqom.pdf}
\end{itemize}

\section{Digital and Audio Resources}
\begin{itemize}
\item AGN: ten Poqomam phrases for the year-end holidays (short phrase list). \url{https://agn.gt/diez-frases-en-pocomam-que-puede-usar-en-las-fiestas-de-fin-de-ano}
\end{itemize}

\section{Curriculum Development Context}
Poqomam is recognised among Guatemala's Mayan languages and the ALMG standardises their orthographies and dictionaries, but little classroom material surfaced. Speaker counts are uncertain (Ethnologue 1991: 27,900 for Southern; the Endangered Languages entry lists 9,548 for 2001). A curriculum would need to start from the ALMG dictionary and ALMG/Ministry of Education bilingual-education (EBI) channels.
\end{document}